
\section{Finite-use graph compilation}
\label{sec:graph}

\subsection{Node inventory}

A practical theorem graph needs only a small family of node constructors.

\begin{center}
\begin{tabularx}{\textwidth}{p{0.20\textwidth}X}
\toprule
Node & Required proof payload\\
\midrule
Exact computation &
Defining equation and preservation of validity/extensionality.\\
Boolean decision &
A reflected quantifier-free formula and Boolean truth-value proof.\\
Bounded search &
A bound, reflected predicate, least-witness specification.\\
Certified iteration &
Initial state, transition equations, QF invariant, clock or residual clock
obligation.\\
Finite relation &
Bounded totality or viability and local compatibility.\\
Guarded operation &
A decidable guard plus proof that the guard implies every side condition.\\
Stable observation &
Pairwise exact, quotient, or error stability on valid witnesses.\\
Residual producer &
A named formula left in the output context with provenance metadata.\\
\bottomrule
\end{tabularx}
\end{center}

This inventory is intentionally close to the common core of the existing
foundation packets.  The difference is that each constructor returns a
proof-carrying module immediately.

\subsection{Adaptive finite use}

Suppose a node queries an input oracle-like interface.  There is no actual
oracle in the object language.  A bounded run stores every query and answer in
its state.  Its transition uses an answer only after checking
\(\ValidObs(q,z)=1\).  The theorem input contains precisely the finite
instances needed along the accepted trace, or one uniform contract from which
those instances are derived by specialization.

If the trace is certificate supplied, the verifier scans its query--answer
pairs.  If the trace is recomputed, the answer producer is a primitive-
recursive function with a proof of its contract.  In either case the final
proof mentions only finitely many specialized facts.

\begin{lemma}[Finite support extraction]
Every concrete \code{Prf PRAAx} derivation has a finite set of hypothesis and
axiom occurrences.  Therefore every compiled theorem instance uses only
finitely many observations of each input interface.
\end{lemma}

\begin{proof}
A derivation is an inductively generated finite tree.  Associate to each
constructor the finite union of the supports of its premises, adding the
formula used by a hypothesis or axiom constructor.  Structural recursion on
the derivation yields a finite list.  The claim is metatheoretic and needs no
new object axiom.
\end{proof}

The significance is operational: an interface may be globally specified, but
the compiler can report the exact finite query support of a generated theorem.

\subsection{The compilation theorem}

Define a well-typed PCFM graph recursively.

\begin{itemize}
\item A leaf is a primitive checked module or a residual requirement.
\item A sequential node applies a local module to compiled predecessors.
\item A branch uses a reflected Boolean condition and two branch modules with
a common conclusion.
\item A loop carries the data of the fixed trace theorem.
\item A bounded existential node invokes least search.
\item A relational node invokes backward viability.
\item An observation node invokes a stability eliminator.
\end{itemize}

\begin{theorem}[Finite-use compilation]
For every finite well-typed PCFM graph \(G\) in free-variable scope \(n\), the
compiler constructs
\[
   \mathsf{compile}(G):
   \Res(G)\vdash_{\PRAAx}\mathsf{Goal}(G),
\]
where \(\Res(G)\) is exactly the list of residual producer formulas not
discharged by a predecessor node.  The construction is stable under term
substitution.  No formula outside the finite support of \(G\) is required.
\end{theorem}

\begin{proof}
Proceed by structural recursion on \(G\).

A primitive module returns its stored derivation.  A residual leaf returns the
hypothesis derivation of its formula.

At a sequential node, recursively compile all predecessors.  The local module
has a context consisting of predecessor outputs and local residuals.  Weaken
the predecessor derivations to the union of residual contexts and apply
\code{Prf.trans}.  Deduplicate residual metadata at the host level; logical
duplication would be harmless.

At a Boolean branch, reflection proves either the condition or its negation
from the Boolean value.  Compile both branches in contexts extended by the
appropriate assumption, then apply intuitionistic disjunction elimination.
If the graph computes the Boolean value explicitly, use finite case analysis
on its Booleanity proof.

At a loop, apply fixed trace reflection from
\cref{sec:trace}, then cut in the terminal contract and clock proof.  The loop
invariant is part of the node's typing data, so compilation cannot proceed
without a quantifier-free proof of the invariant and its base and step
derivations.

At a bounded existential node, invoke bounded least witness
(\cref{sec:boolean}) and specialize the continuation at the returned witness.
At a relational node use the viability theorem of
\cref{sec:relational}.  At an observation node use exact, quotient, or
error-stable erasure.

All recursive calls return actual derivations, and every assembly step is a
constructor or checked structural theorem of \code{Prf}.  Thus the result has
the displayed type.  Substitution stability follows by applying
\code{Prf.subst} to the compiled result; alternatively it can be proved
constructor by constructor, using substitution compatibility of Boolean
reflection, traces, and search.  Finiteness of support follows from finiteness
of the graph and the preceding lemma.
\end{proof}

\subsection{No hidden semantic trust in the compiler}

The compiler itself is ordinary Lean code and may contain bugs.  A scheduling
bug can choose the wrong predecessor; a demand analysis bug can compute an
insufficient precision; a serialization bug can parse a malformed graph.
None of these bugs can produce a theorem of the wrong type unless the Lean
kernel accepts a \code{Prf PRAAx} term.

This is the main advantage of the shallow design.  The compiler is not in the
trusted base for mathematical soundness.  The trusted base consists of Lean's
kernel, the definitions of \code{Prf} and \code{PRAAx}, and the imported
checked derivations.  Compiler correctness remains important for completeness,
performance, and whether the intended theorem was formalized, but not for the
validity of an accepted proof term.

\subsection{Residual minimality}

The phrase ``exactly the residual producer formulas'' should be understood
operationally.  The compiler removes a requirement when a predecessor
derivation is supplied.  It does not decide logical redundancy among the
remaining formulas.  A later simplifier may prove that one residual follows
from others and discharge it by another module.

This choice avoids an undecidable minimization problem.  It also preserves
provenance: a theorem's ledger records the boundaries the graph genuinely
used, even if some are mathematically redundant.

\subsection{Caching and theorem reuse}

A compiled module may be cached by:
\[
  (\text{repository commit},\text{module name},
    \text{normalized input syntax},\text{proof hash}).
\]
Caching proof terms is an engineering optimization.  On reuse, Lean still
checks the imported theorem.  Frequently used nodes---finite folds, rational
inequalities, matrix operations, power-series truncations, interval
refinements---should live in \code{IPRADerivations}, while theorem-specific
graphs remain near their packets.

\subsection{From graphs to the native scalar calculus}

Once a result has been proved as
\(\Gamma\vdash_{\PRAAx}\varphi\), the checked theorem
\code{pra_iff_native} translates it to the native scalar calculus.  PCFM
therefore need not duplicate every constructor for the smaller calculus.
The ordinary interface is better for development; the native endpoint remains
available for proof-checker size experiments, Busy Beaver exports, or a final
minimal certificate.
